% !TeX root = ../main.tex

\chapter{引言}
\section{研究背景与问题}
以合理价格购买持续成长企业，是价值与成长相结合的选股思路。柏顿·墨基尔成功选股法在A股的量化实现，以历史盈余成长、当前成长排名和相对低市盈率共同限制候选股票。申万宏源研究[1]提供了可复现的规则与回溯收益，适合检验投资理念在本地数据与具体程序中的表现。

复现的难点不只在于得到一条上涨净值曲线。历史记录是否齐全、百分比排名以何种股票为分母、负PE是否参与排序、候选稀缺时怎样构造组合，都会改变最终名单。在此基础上引入多因子得分和动态仓位，还需要区分选股改变与风险暴露调整，否则容易把多项同时发生的变化归因于某一个“有效因子”。

本文首先回答三个问题：第一，基础规则在原复现区间与延展区间的收益和风险是否一致；第二，候选池内排名与全截面独立排名如何改变组合；第三，在两种选股路径上叠加同一仓位规则，能否以可接受的收益代价降低风险。研究同时保留不利结果，不以更高回溯收益作为判断规则是否正确的标准。

\section{文献联系与本文定位}
本文直接复现对象为券商研究[1]。Fama和French[2]的横截面收益研究提示，组合的高收益还需放在市场和风格暴露背景下理解；价值与动量结合的研究[3]为多维选股提供相关背景，但不能直接证明本项目质量、动量与规模等权得分的有效性。本文没有将其他市场的经验结果移植为A股结论，也未额外引入新的定价模型。新增差异核验以原文、数据覆盖及独立计算作为证据，不根据收益率接近程度决定规则。

本文的工作在于建立可核查的规则链条。基础复现与延展保持同一筛选逻辑；路径一在已修正正PE与初始期规则的代码上保留19/20和三因子选择；路径二仅改变两项排名及各自分母；路径三、四分别只改变前两条路径的仓位。成对比较可以明确实现差异的经济幅度，但不等于随机实验或单因子的因果识别。

\section{原研报与基础复现的基本情况对照}
本研究的直接复现对象是申万宏源2015年发布的柏顿·墨基尔成功选股法[1]，比较对象首先应是基础复现B。延展E改变了结束时间，P1—P4还改变了部分选股或仓位规则，均不能直接替代B来判断原策略是否复现。表1列出研究设计，表2列出原文公布值与本项目实际计算值；原文未披露的细节按“未详”保留，不作一致性假定。

\begin{table}[htbp]
\centering
\caption{申万原研报与基础复现的研究设计对照}
\renewcommand{\arraystretch}{1.18}
\small
\setlength{\tabcolsep}{4.5pt}
\begin{tabularx}{\textwidth}{@{}p{2.45cm}LL@{}}
\toprule
比较项目 & 申万原研报披露 & 本研究基础复现B \\
\midrule
研究对象 & 历史全部A股，含后来退市股票 & 历史A股范围；按现有CSMAR估值、财务和收益覆盖形成可计算集合 \\
成长率 & 净利润同比增长率 & 累计总净利润同比；上年同期利润非正或缺失时，不形成有效成长率 \\
估值指标 & 股价/每股收益（TTM） & 月末总市值/总净利润（TTM）；归属与股本口径尚未逐项确认等价 \\
成熟期筛选 & 过去五年持续正成长，成长率前30\%；正PE且低估值30\% & 最近20季齐全且全为正；当期成长率前30\%与估值最低30\%取交集 \\
初始五年 & 历史每期成长率为正，成长率最高前15只 & 先取历史合格与估值条件交集，再按成长率选最多15只 \\
排名与缺失值 & 市场相对排序；负PE是否占名次、缺失历史处理等细节未完全披露 & 成长和PE各自全截面排名；负PE先占位，再要求入选PE为正；缺失不视作正成长 \\
调仓时点 & 每年四次；全市场最晚财报发布日与报告期后90个交易日孰早 & 固定3/6/9/12月末映射当季财报，次月起持有三个月 \\
权重与费用 & 资金平均配置；交易费率0.13\%，逐笔实现未详 & 每月对有收益记录的持仓等权；有持仓季度首月固定扣0.13\% \\
收益与基准 & 日频回测；上证综指 & 个股含红利月收益；上证综指日收益复利为月收益 \\
\bottomrule
\end{tabularx}
\par\smallskip
\begin{minipage}{\textwidth}\small 资料来源：申万原研报[1]第5—6页及本项目基础程序。初始五年的先后筛选顺序、亏损PE的排名占位和历史资料覆盖，原文不能唯一确定；本表如实列出本项目实现。报告期末财务并非当日已公开信息，本研究调仓近似不等价于原文发布日规则。\end{minipage}
\end{table}

\begin{table}[htbp]
\centering
\caption{申万原研报与基础复现的主要结果对照}
\renewcommand{\arraystretch}{1.18}
\small
\setlength{\tabcolsep}{4.5pt}
\begin{tabularx}{\textwidth}{@{}p{2.05cm}LLL@{}}
\toprule
比较项目 & 申万原研报披露 & 本研究基础复现B & 比较说明 \\
\midrule
实际统计区间 & 原文称2003年至2015-07-31；2938个交易日 & 2003-04至2015-06；147个月 & 起点细节及终点未完全对齐 \\
年化收益 & 36.74\% & 38.86\% & 相差2.12个百分点；不是同日同频差额 \\
累计收益 & 4176.87\% & 5481.49\% & 须同时阅读统计区间 \\
最大回撤幅度 & 70.19\% & 70.11\% & 原日频、本月末；相近不证明路径相同 \\
夏普比率 & 0.93 & 1.100（零无风险利率） & 频率及无风险口径未统一 \\
平均持股数 & 7只 & 8.02只（49次信号等权平均） & 原文均值汇总细节未披露 \\
最少/最多持股 & 0/15只 & 0/15只 & 仅比较观测极值，非容量规则相同 \\
空仓情况 & 2段空仓，具体起止见下文 & 3个空仓季度信号，共9个月 & 日历时段和信号数不属同一单位 \\
末期组合 & 2015-05-25调仓；2015-07-31持有8只 & 2015-03-31信号；持有9只至6月 & 共有4只；选股日和观察终点不同 \\
\bottomrule
\end{tabularx}
\par\smallskip
\begin{minipage}{\textwidth}\small 资料来源：原研报[1]第6、8、12页；本研究六方案月收益、逐季持股表及统一指标。最大回撤为便于对照取损失幅度正值，后文绩效表以负值表示回撤。B的CAGR使用147/12年；原文的精确建仓日、年化日数约定及夏普无风险利率未充分披露，原图首个刻度不作建仓日依据。\end{minipage}
\end{table}

收益方面，B的年化38.86\%与原文36.74\%在数量级上接近，但累计收益、调仓日期和末期名单均未完全一致，不能据此宣布精确复现。风险方面，约70\%的峰谷损失是两者共同呈现的显著特征；日频与月末回撤相近只表明本次数值接近，不能证明月内路径或风险控制能力相同。原文夏普0.93与本研究零利率月频夏普1.100也不作严格优劣判断。

持仓方面，原文两段空仓为2009年11月2日至2010年6月29日、2010年9月1日至10月29日；B的空仓信号为2009年9月30日、2010年3月31日、2013年3月29日，对应2009年第四季度、2010年第二季度及2013年第二季度各三个月。两者有部分时间重叠，但空仓安排明显不同。末期共有的4只为中国医药（600056）、南京高科（600064）、欧亚集团（600697）和中国平安（601318）；研报另有康力电梯、海康威视、老板电器和奥瑞金，当前资料下未通过本项目历史成长条件。第8章进一步区分已确认的规则差异、已定位的历史覆盖缺口与仍未解决的时点问题。

因此，本文将B定位为“保留原选股思想、明确月频近似的规则复现”，将E作为同规则历史延长。原研报36.74\%是其原样本与实施口径下的公布结果；本研究E的26.17\%另含2015年后的市场经历，不能据两者差额认定延展失败，也不能把B较高的历史年化解释为优于原模型。
